\documentclass[11pt,a4paper,headinclude,footinclude,
%DIV16,
normalheadings]{scrreprt}

\usepackage[automark]{scrpage2}
\usepackage[ansinew]{inputenc}
%\usepackage{german}
%\usepackage{bibgerm}
\usepackage{amsmath}
\usepackage{amsfonts}
\usepackage{theorem}
\usepackage{color}
\usepackage{listings}
 \lstset{language=C++, basicstyle=\ttfamily\scriptsize,
   keywordstyle=\color{black}\bfseries, tabsize=4,
   stringstyle=\ttfamily, commentstyle=\it, extendedchars=true,
   breaklines=true,escapeinside={/*@}{@*/}
   }
\usepackage{hyperref}
\usepackage{breakurl}
\usepackage{psfrag}
%\usepackage{makeidx}
\usepackage{graphicx}
\DeclareGraphicsExtensions{.eps, .jpg}

\setlength{\parskip}{1.5ex plus0.5ex minus0.5ex}   % Abstand zwischen Abs�tzen
\setlength{\parindent}{0pt}                        % keine Einr�ckung der ersten Absatzzeile

% shrink boxes to vertical line
\def\clap#1{\hbox to 0pt{\hss#1\hss}}
\def\mathclap{\mathpalette\mathclapinternal}
\def\mathclapinternal#1#2{%
    \clap{$\mathsurround=0pt#1{#2}$}}

% fuer Zitate
\usepackage[round]{natbib}

\newcommand{\C}{\mathbb{C}}
\newcommand{\R}{\mathbb{R}}
\newcommand{\N}{\mathbb{N}}
\newcommand{\Z}{\mathbb{Z}}
\newcommand{\Q}{\mathbb{Q}}
\newcommand{\Th}{\mathcal{T_h}}
\newcommand{\Dune}{\texttt {DUNE}}
\newcommand{\Grid}{\texttt {dune-grid}}
\newcommand{\Fem}{\texttt {dune-fem}}
\newcommand{\todo}{{\color{red}{\quad !Here is still something to do!}}}

%The theorems
\theorembodyfont{\upshape}
\theoremheaderfont{\sffamily\bfseries}
\newtheorem{exc}{Excercise}
\newtheorem{rem}{Remark}
\newtheorem{lst}{Listing}
\newtheorem{warn}{Warning}

\pagestyle{scrheadings}

%%%%%%%%%%%%%%%%%%%%%%%%%%%%%%%%%%%%%%%%%%%%%%%%%%%%%%%%%%%%%%%%

% titlepage

%%%%%%%%%%%%%%%%%%%%%%%%%%%%%%%%%%%%%%%%%%%%%%%%%%%%%%%%%%%%%%%%


\title{Introduction to \Fem\\
\vspace{10mm}
\includegraphics{pics/dune-logo1}}

\author{The Freiburg \Dune\ Group$^\ast$}

\date{\today}

\publishers{
{\normalsize $^\ast$Abteilung f�r Angewandte Mathematik, Universit�t Freiburg,\\
Hermann-Herder-Str.~10, D-79104 Freiburg, Germany}\\
%
\bigskip
{\normalsize \url{http://www.dune-project.org/}}
}


%\makeindex

%%%%%%%%%%%%%%%%%%%%%%%%%%%%%%%%%%%%%%%%%%%%%%%%%%%%%%%%%%%%%%%%

% \begin{document}

%%%%%%%%%%%%%%%%%%%%%%%%%%%%%%%%%%%%%%%%%%%%%%%%%%%%%%%%%%%%%%%%

\begin{document}
  \maketitle

\begin{abstract}
Dieses Dokument stellt in der vorliegenden Form keine vollst�ndige Referenz in \Fem\ dar, sondern es soll als eine Art roter Faden dienen, damit \Dune-Neulinge eine Idee bekommen, was alles m�glich ist und wo man anfangen kann. Das Problem aus meiner Sicht ist, dass die vorhandene Dokumentation auf zu viele verschiedene Orte verteilt und teilweise nur schwer zu finden ist. Dementsprechend fasst dieses Dokument zu einem grossen Teil einfach Informationen aus verschiedenen Quellen zusammen oder verweist auf andere Quellen.

Der erste Teil von Kap.1 ist fast w�rtlich von der \Dune\ Homepage kopiert. Kapitel 2 verwendet die Programmbeispiele aus der Dune Summer School 2007 von R. Kl�fkorn, und Kapitel 3 ebenso.

Worauf also Wert gelegt wurde, sind einfache Beispiele und viele Verweise auf weitere Dokumentation und Hilfedateien, insbesondere also auf die \texttt{doxygen}-Dokumentation.

\begin{flushright}Freiburg, Januar 2008\end{flushright}
\end{abstract}


  \tableofcontents

  %Chapter 1
  \chapter{What is \Dune?}
  \Dune, the Distributed and Unified Numerics Environment is a modular toolbox for solving partial differential equations with grid-based methods. It supports easy discretization using methods like Finite Elements, Finite Volumes, and also Finite Differences.

\Dune\ is free software licensed under the GPL with a so called "runtime exception" (\url{http://www.dune-project.org/license.html}).

The main intention is to create slim interfaces allowing an efficient use of legacy and/or new libraries. Using C++ techniques \Dune\ allows to use very different implementations of the same concept (i.e. grids, solvers, ...) using a common interface with a very low overhead. More: \url{http://www.dune-project.org/dune.html} \cite{duneweb}.



\section{Core Modules}
The framework consists of a number of modules which are downloadable as separate packages. The current core modules are:

\begin{itemize}
 \item \texttt{dune-common}\\
    contains the basic classes used by all \Dune-modules. It provides some infrastructural classes for debugging and exception
    handling as well as a library to handle dense matrices and vectors.
 \item \Grid\\
    is the most mature module. It defines nonconforming, hierarchically nested, multi-element-type, parallel grids in
    arbitrary space dimensions. Graphical output with several packages is available, e.g. file output to IBM data explorer
    and VTK (parallel XML format for unstructured grids). The graphics package Grape has been integrated in interactive mode.
 \item \texttt{dune-istl} (Iterative Solver Template Library)\\
    provides generic sparse matrix/vector classes and a variety of solvers based on these classes. A special feature is the use
    of templates to exploit the recursive block structure of finite element matrices at compile time. Available solvers include
    Krylov methods, (block-) incomplete decompositions and aggregation-based algebraic multigrid.
\end{itemize}




\section{Available Grid Implementations}\label{availablegrids}
 So far six grid implementations are available through the \Grid interface. Each is geared towards a different purpose.
\begin{itemize}
    \item \texttt{ONEDGRID}: A sequential locally adaptive grid in one space dimension
    \item \texttt{SGRID}: A structured grid in n space dimensions
    \item \texttt{YASPGRID}: A structured parallel grid in n space dimensions (used as default grid)
    \item \texttt{ALUGRID\_CUBE, ALUGRID\_SIMPLEX}: A parallel locally adaptive grid in 2 and 3 space dimensions
    \item \texttt{ALBERTAGRID}: The grid manager of the Alberta toolbox
    \item \texttt{UGGRID}: The grid manager of the UG toolbox (UG is not freely available)
\end{itemize}
More information about these grids: \url{http://www.dune-project.org/doc/devel/features.html}.



% This file is included by
%    -  1_what_is_dune.tex and therefore by femhowto.tex
%    -  installation-fem.tex to provide a nice installation guide for the internal website

\section{Download and installation}
This section describes how members of the Section of Applied Mathematics in Freiburg can create a local working installation of \Dune\ .

These are only the basic steps, maybe you have to tune up some options. For non-Freiburg users, the steps are similiar.
For more information, see \url{http://www.dune-project.org/doc/installation-notes.html}.

\begin{itemize}
\item Create a directory for your \Dune\ installation and change to it. For example:
\begin{lstlisting}
mkdir dune
cd dune
\end{lstlisting}

% \item Download the latest (stable) core modules from the official \Dune\ homepage: \url{http://www.dune-project.org/download.html}.
%    The module DUNE ISTL is usually not needed for Freiburg.
%
% \item Copy the tar.gz Archives to your \Dune\ directory and extract them. For example:
% \begin{lstlisting}
% tar -xzf dune-common-1.0.tar.gz
% tar -xzf dune-grid-1.0.tar.gz
% tar -xzf dune-grid-howto-1.0.tar.gz
% \end{lstlisting}
%
% \item Rename the resulting directories:
% \begin{lstlisting}
% mv dune-common-1.0 dune-common
% mv dune-grid-1.0 dune-grid
% mv dune-grid-howto-1.0 dune-grid-howto
% \end{lstlisting}
%
% \item Download \Dune\ \Fem\ via svn: Just type in your \Dune\ directory:
% \begin{lstlisting}[language=csh]
% svn checkout svn+ssh://morgoth/raid5/dune/src/svn-archive/
%       dune-fem/release-0.9 dune-fem
% svn checkout svn+ssh://morgoth/raid5/dune/src/svn-archive/dune-femhowto
% \end{lstlisting}

\item Checkout the necessary \Dune\ repositories via \lstinline!svn!. To do this, type the following commands in your \Dune\ directory:
  \begin{lst}[File ./installation/checkout.sh] \label{checkout.sh} \mbox{}
    \lstinputlisting[language=bash]{./installation/checkout.sh}
  \end{lst}
Alternatively, you can download the corresponding tar.gz archives from \url{http://www.dune-project.org/download.html} and \url{http://dune.mathematik.uni-freiburg.de} and extract them into your \Dune\ directory.


\item Create a file named \lstinline!config.opts! with the following content in your \Dune\ directory. Consider this file as an example, that will work correctly only for members of the Section of Applied Mathematics in Freiburg! Other users have to adapt (or to comment out) at least the paths to the external modules.

  \begin{lst}[File ./installation/config.opts] \label{config.opts} \mbox{}
    \lstinputlisting{./installation/config.opts}
  \end{lst}

\item Finally, configure und compile \Dune\ using the \lstinline!dunecontrol! script: Type in your \Dune\ directory:
\begin{lstlisting}
       ./dune-common/bin/dunecontrol --opts=config.opts all
\end{lstlisting}

   The script needs some minutes to finish. After finishing, you can start working!
\end{itemize}


You can find the file \texttt{config.opts} and shell scripts with the commands described above under
\url{http://www.mathematik.uni-freiburg.de/IAM/Research/projectskr/dune/freiburg_intern/tools.html}
or in the subdirectory \texttt{doc/installation} of the \texttt{dune-femhowto} module.




\section{Create your own project}
You can create your own \Dune\ project by using the \lstinline!duneproject! script. Type in your \Dune\ directory
\begin{lstlisting}
           ./dune-common/bin/duneproject
\end{lstlisting}
and follow the instructions. After creating your project, you have to rerun the \lstinline!dunecontrol! script:
\begin{lstlisting}
       ./dune-common/bin/dunecontrol --opts=config.opts --module=YOUR_MODULE_NAME all
\end{lstlisting}
Take a look in your new \Dune\ project directory. The \lstinline!duneproject! script created a sample source code file, which was compiled by \lstinline!dunecontrol!.

If you want to write your own code, replace the sample source code with your own code, and modify the file
\lstinline!Makefile.am!, if necessary. You can also create your own subdirectories, but then you additionally have to
modify the file \lstinline!configure.ac! and run again \lstinline!dunecontrol! in order to create the makefiles in the subdirectories.
If you are unsure how to do these modifications, look how this is done in the core \Dune\ Modules,
or read the \Dune\ Build System Howto on \url{http://www.dune-project.org/doc/buildsystem/buildsystem.pdf}.






  %Chapter 2
  \chapter{Short Introduction to the \Grid\ Interface}
  This chapter gives a really short introduction to \Grid. For a more detailed explanation please read the \texttt{DUNE GRID HOWTO}
(\url{http://www.dune-project.org/doc/grid-howto/grid-howto.pdf}).

Understanding code can be sometimes difficult. The \Dune\ Coding Style document can help you to understand \Dune\ code, see 
\url{http://www.dune-project.org/doc/devel/codingstyle.html}.

The source code to all of the following examples is shipped with this documentation, see directory \lstinline!dune-femhowto/src_grid!.


\section{The \texttt{dune hello world} programm}
  \begin{lst}[File ../src\_grid/dune\_hello\_world.cc] \label{lsthelloworld} \mbox{}
    \lstinputlisting[numbers=left,numberstyle=\tiny,numbersep=5pt]{../src_grid/dune_hello_world.cc}
  \end{lst}

This little programm has a sequentiell and a parallel part. The parallel part is regarded in section \ref{parallel}. 
In fact, this program just do nothing. But you can check whether your \Dune\ installation is working. Output:
\begin{lstlisting}
Hello World! This is DUNE.
This is a sequential program.
\end{lstlisting}




\section{Getting started with your first grid}
Now we want to create our first grid:

  \begin{lst}[File ../src\_grid/gettingstarted.cc] \label{lstgettingstarted} \mbox{}
    \lstinputlisting[numbers=left,numberstyle=\tiny,numbersep=5pt]{../src_grid/gettingstarted.cc}
  \end{lst}

Here we create an SGRID (lines 17-21) for several space dimensions (lines 50,53 and 56) und print some information about these grids. 
Notice we use generic programming:
\begin{lstlisting}
template <int dim> 
struct Info 
   { ... };
\end{lstlisting}

Interesting is the number of codimensions in the several space dimensions. We just created an unitcube for the 
dimension 1, 2 and 3. 
\begin{itemize}
 \item 1D: We have one codim 0 element (the unit intervall) and two codim 1 elements (the end points of the intervall)
    \begin{lstlisting}
      => SGrid(dim=1,dimworld=1)
      level 0 codim[0]=1 codim[1]=2
      leaf    codim[0]=1 codim[1]=2
    \end{lstlisting}
 \item 2D: We have one codim 0 element (the unit square), four codim 1 elements (edges) and four codim 2 elements (vertices).
    \begin{lstlisting}
      => SGrid(dim=2,dimworld=2)
      level 0 codim[0]=1 codim[1]=4 codim[2]=4
      leaf    codim[0]=1 codim[1]=4 codim[2]=4
    \end{lstlisting}
 \item 3D: We have one codim 0 element (the unit cube), six codim 1 elements (faces), twelve codim 2 elements (edges) 
           and eight codim 3 elements (vertices).
    \begin{lstlisting}
      => SGrid(dim=3,dimworld=3)
      level 0 codim[0]=1 codim[1]=6 codim[2]=12 codim[3]=8
      leaf    codim[0]=1 codim[1]=6 codim[2]=12 codim[3]=8
    \end{lstlisting}
\end{itemize}
Because we did't refined the grid, there only exists a level 0 grid, and this level grid is the same as the leaf grid.
You have access to all these different codim elements via iterators, even after refining you have access to all elements of the different levels.

For a more detailed explanation please read the \Dune\ GRID HOWTO 
(\url{http://www.dune-project.org/doc/grid-howto/grid-howto.pdf}) 



\section{The DGF-Parser: Using different grids / macrogrids}
\Dune\ is designed to easily handle with different grid implementations. Therefore it is really easy to change the used grid.
To show this fact we take a look to the following code:

  \begin{lst}[File ../src\_grid/dgfparser.cc] \mbox{}
    \lstinputlisting[numbers=left,numberstyle=\tiny,numbersep=5pt]{../src_grid/dgfparser.cc}
  \end{lst}

By default, programs are compiled using the options \lstinline!GRIDTYPE=YASPGRID GRIDDIM=3!. If you want to change the grid type or the number of space dimensions, you just have to recompile the program again using some other options for make, for example: \lstinline!make clean GRIDTYPE=ALUGRID_SIMPLEX GRIDDIM=3!. How you have to implement this feature is shown in \lstinline!dgfGridType()!, lines 45-56.
Attention: For changing your grid in the described way, you must define \lstinline!GRIDTYPE! somewhere in an \lstinline!makefile.am! or in your configuration options during installation!

Another possibility is shown in \lstinline!dgfTest()!, lines 22-30. Here the used grid is hardcoded in the source code (line 29). This is similar to the version in Listing \ref{lstgettingstarted}. The only difference is that you use there an SGrid without macrogrids and without the DGF-Parser.

For more information: See DGF-Parser documentation (Link is below). We also refer to the list of the available grid implementations (\ref{availablegrids}).

Are you bored about these stupid unit cubes? You can change your macrogrid just by editing the corresponding \texttt{dgf} files. How this can be done is explained in the documentation of the DGF-Parser. This documentation you find in your local \Grid\ documentation, under the point \texttt{I/O - The Dune Grid Format (DGF)}, or online under \url{http://www.dune-project.org/doc/doxygen/dune-grid-html/group__DuneGridFormatParser.html}. 

You can find the example \texttt{dgf}-files from that documentation in your local \Fem\ installation, in the subdirectory \texttt{macrogrids/DGFMacrogrids}.

The program above also contains an output routine, which writes the grid information into a \texttt{vtk} file (lines 37-38 and 69-70). You can visualize your grid by opening these files with \texttt{paraview} (\url{http://www.paraview.org/}) or another VTK-Viewer.

\begin{exc}
Play with different dimensions, different grids implementations and different macrogrids, and visualize the results with \texttt{paraview}!
\end{exc}


\section{The parallel \texttt{dune\_hello\_world} programm}\label{parallel}
Now we want to run the \texttt{dune\_hello\_world} (listing \ref{lsthelloworld}) programm parallel on several computers.
\bigskip

To do that, you have to recompile the package with a modified \texttt{config.opts}. Replace in listing \ref{config.opts} \lstinline!--disable-parallel! with \lstinline!--enable-parallel! (line 20) and \lstinline!CXX=g++! with \lstinline!CXX=mpiCC! (line 27). After that, recompile the \Fem\ HOWTO with your modified \texttt{config.opts}, which ist now called \texttt{config\_parallel.opts}:
\begin{lstlisting}
  ./dune-common/bin/dunecontrol --opts=config_parallel.opts --only=dune-femhowto all
\end{lstlisting}
Notice: By using the \texttt{--only=XXXX} option, dunecontrol will only work on the module XXXX. 


\begin{enumerate}

\item Running the program on one computer:
\begin{lstlisting}
./dune_hello_world
\end{lstlisting}

Output:
\begin{lstlisting}
Hello World! This is DUNE.
I am rank 0 of 1 processes!
\end{lstlisting}
Compare that with the output of the serial version!

\item Running on several computers using the ''Simple Linux Utility for Resource Managment (SLURM)'':
\begin{lstlisting}
srun -A -n 4
mpirun nice ./dune_hello_world
exit
\end{lstlisting}
Here we use 4 machines (\lstinline!-n 4!) for running the program, resulting in an output similar to this:
\begin{lstlisting}
Hello World! This is DUNE.
Hello World! This is DUNE.
I am rank 0 of 5 processes!
Hello World! This is DUNE.
I am rank 3 of 5 processes!
I am rank 1 of 5 processes!
Hello World! This is DUNE.
I am rank 2 of 5 processes!
Hello World! This is DUNE.
I am rank 4 of 5 processes!
\end{lstlisting}
Don't forget the final \lstinline!exit! to terminate your job!!!!

\item Running on several computer without using SLURM:
\begin{lstlisting}
nice mpirun -np 4 ./dune_hello_world
\end{lstlisting}
The output will be the same as above. Using SLURM has some advantages, for example it can avoid conflicts when running several jobs on one machine. So better use it! More information: SLURM manpages (type: \lstinline!man slurm!).
\end{enumerate}

You should also use \lstinline!nice! like in the examples above to adjust scheduling priority.

Some maybe helpful commands:
\begin{itemize}
 \item \lstinline!squeue!: Returns the list of jobs managed by SLURM (and their job IDs).
 \item \lstinline!scancel <JOB ID>!: Cancels the corresponding job.
 \item \lstinline!sinfo!: Print some general information.
\end{itemize}
\bigskip

Attention: Only YaspGrid and AluGrid are able to perform parallel computations!

For real applications please read section 8 in the \Dune\ GRID HOWTO   (\url{http://www.dune-project.org/doc/grid-howto/grid-howto.pdf}).

More information: \texttt{dune-common} documentation - \texttt{Parallel Communication}  (\url{http://www.dune-project.org/doc/doxygen/dune-common-html/group__ParallelCommunication.html}).


\section{Implementation of a Finite Volume Scheme with \Grid}
You can find different implementations of a Finite Volume Scheme in the \Dune\ GRID HOWTO 
(\url{http://www.dune-project.org/doc/grid-howto/grid-howto.pdf}).

\begin{itemize}
 \item Basic scheme: Section 6.3
 \item Adaptive scheme: Section 7.2
 \item Parallel scheme: Section 8.3
\end{itemize}










  %Chapter 3
  \chapter{The Poisson-Example with \Fem}
  %
% Written by K. Michel
%

In this chapter we want to show how to use \Fem\ for solving a simple test problem.
As a test problem, we choose the Poisson problem: \label{poissonproblem}

\begin{eqnarray}
  -\triangle u &=& f  \quad \mbox{in}\ \Omega\\
             u &=& g  \quad \mbox{on}\ \partial\Omega
\end{eqnarray}
We want to solve this problem with the finite element method using Lagrangian elements. The implementation of this method is done in the files in the directory \lstinline!dune-femhowto/src_poisson!.

The numerical treatment of this problem is described in chapter 1 of this skript (in german): \url{http://www.mathematik.uni-freiburg.de/IAM/Teaching/ubungen/sci_com_SS06/skriptum.pdf}.

In this example, we choose the following problem data:
\begin{eqnarray}
  \Omega &:=& ]0,1[^{dimworld}, \quad x = (x_1, ..., x_{dimworld}) \\
  f(x) &:=& 2 \sum_{i=1}^{dimworld} \prod_{j \neq i} (x_j-x_j^2)  \quad \forall x \in \Omega  \\
  u(x) &:=& \prod_{i=1}^{dimworld} (x_i - x_i^2)  \quad \forall x \in \Omega   \\
  g(x) &:=& u(x) \quad \forall x \in \partial\Omega
\end{eqnarray}
where $u$ is the exact solution. The boundary values are just the values of the exact solution on the boundary points.

Note: In the current implementation we implemented another right hand side $f$.
\begin{eqnarray}
  f(x,y) &:=& 8 \pi^2 \cos(2 \pi x) \cos(2 \pi y)  \quad \forall x \in \Omega  \\
  u(x,y) &:=&         \cos(2 \pi x) \cos(2 \pi y)  \quad \forall x \in \Omega   \\
  g(x,y) &:=& u(x,y) \quad \forall x \in \partial\Omega
\end{eqnarray}
The second problem data should better be used only with \texttt{dimworld=2}. For changing between these two versions, please modify \texttt{problem.hh}.



\section{Implementing the algorithm and the problem data}

\subsection{The algorithm in \texttt{lagrange.cc} }
\subsubsection{The function \texttt{main()} }
We want to skip all the \texttt{includes} at the beginning of the file and jump directly to the \texttt{main} routine.

\begin{itemize}
 \item Initialize grid and some more initialization stuff

 \item Call \texttt{Algorithm::calc()} twice for EOC calculation:
\begin{lstlisting}
for( int i=0; i<2; ++i )
{
	// calculate L2-projection
	projectionError[i] =
       Algorithm<GridType, L2ProjectionProblem, MassOperator>
          ::calc(grid,"l2-projection", false);

	// calculate poisson problem
	poissonError[i] =
       Algorithm<GridType, PoissonProblem, LaplaceOperator>
          ::calc(grid,"poisson",true);

	// make refinement for next step
	grid.globalRefine( DGFGridInfo<GridType>::refineStepsForHalf() );
}
\end{lstlisting}

 \item Calculate EOC:
\begin{lstlisting}
const double poissonEOC =
                log( poissonError[0] / poissonError[1] ) / M_LN2;
const double projectionEOC =
                log( projectionError[0] / projectionError[1] ) / M_LN2;
\end{lstlisting}

\end{itemize}
\texttt{L2ProjectionProblem} and \texttt{PoissonProblem} are defined in \texttt{problem.hh},
\texttt{MassOperator} and \texttt{LaplaceOperator} are defined in \texttt{massmatrix.hh} respective \texttt{laplace.hh}. These files will be explained in detail later.


\subsubsection{The struct \texttt{Algorithm} }
\begin{lstlisting}
template < class GridImp,
           template <class> class ProblemImp,
           template <class> class OperatorImp >
struct Algorithm
\end{lstlisting}
This struct contains some (important) typedefs and the method \texttt{calc()}, which holds the algorithm.
Explanation of the typedefs:

\begin{itemize}
 \item \lstinline!typedef GridImp GridType!\\
      New name for template parameter.

 \item \lstinline!typedef LeafGridPart<GridType> GridPartType!\\
      A \texttt{GridPart} is a subset of entities of the whole grid. Here our \texttt{GridPart} is build up by all leaf elements.

 \item \lstinline!FunctionSpace<double, double, GridType::dimension, 1>  FunctionSpaceType!\\
      The analytical functionspace we want to use: We use functions $f: \R^{dimension} \rightarrow \R $,
      where Elements from $\R$ are represented by \texttt{doubles}.

 \item \lstinline!typedef LagrangeDiscreteFunctionSpace<FunctionSpaceType, GridPartType, POLORDER>!\\
       \lstinline!    DiscreteFunctionSpaceType!\\
      The discrete function space we want to use, here the lagrange space.

 \item \lstinline!typedef AdaptiveDiscreteFunction<DiscreteFunctionSpaceType> DiscreteFunctionType!\\
      Build up discrete functions from the discrete function space.

 \item \lstinline!typedef ProblemImp<FunctionSpaceType> ProblemType!\\
       \lstinline!typedef typename ProblemType::RHSFunctionType RHSFunctionType!\\
       \lstinline!typedef typename ProblemType::ExactSolutionType ExactSolutionType!\\
      The problem data and data functions (as defined in problem.hh)

 \item \lstinline!typedef OperatorImp<DiscreteFunctionType> OperatorType!\\
      New name for template parameter (will be the laplace or the massmatrix operator).

 \item \lstinline!typedef OEMBICGSTABOp<DiscreteFunctionType,OperatorType> InverseOperatorType!\\
      Define the type of inverse operator we are using to solve the system. There are different inverse operators available.
\end{itemize}


\subsubsection*{ Explanation of the algorithm in the method \texttt{calc()} }
\begin{lstlisting}
static double calc (GridType& grid,
                    const std::string name,
                    bool verbose = false )
\end{lstlisting}

\begin{itemize}
 \item Create instances of the types defined above:
	\begin{itemize}
	  \item \texttt{discreteFunctionSpace} of type \texttt{DiscreteFunctionSpaceType} (and print also some information)
	  \item \texttt{solution} and \texttt{rhs} of type \texttt{DiscreteFunctionType} (discrete functions for solution and right hand side)
	  \item \texttt{f} of type \texttt{RHSFunctionType} (analytical right hand side)
	\end{itemize}

 \item Build right hand side (see boundary.hh):
\begin{lstlisting}
assembleRHS<2*DiscreteFunctionSpaceType::polynomialOrder> ( f , rhs );
\end{lstlisting}

 \item Set the dirichlet boundary points to the corresponding values of the exact solution $u$ by modifying \texttt{rhs}
       (see boundary.hh, problem.hh):
\begin{lstlisting}
ExactSolutionType u( discreteFunctionSpace );
// ...
for( IteratorType it = discreteFunctionSpace.begin(); it != endit; ++it )
{
  boundaryTreatment( *it, u, rhs );
}
\end{lstlisting}

 \item Create instance of our problem operator (here: Laplace or MassMatrix Operator),
       see \texttt{laplace.hh} and \texttt{massmatrix.hh}.
\begin{lstlisting}
OperatorType problemOperator( discreteFunctionSpace );
\end{lstlisting}

 \item Solve the linear system by using the inverse operator (defined above) of our \texttt{problemOperator}:
\begin{lstlisting}
InverseOperatorType cg( problemOperator, 1e-8, 1e-8, 20000, verbose );
cg( rhs, solution );
\end{lstlisting}
For more information see the doxygen documentation of \Fem\ (Operators - Linear solver)

 \item Error calculation and graphical output.
\end{itemize}




\subsection{The problem data in \texttt{problem.hh} }
First we implement the problem data: The right hand side $f$ and the exact solution $u$ of our Poisson problem (see page \pageref{poissonproblem}). The exact solution is needed for the EOC calculation and also for setting the boundary values.
\begin{lstlisting}
template <class FunctionSpaceType>
class myRHSFunction
    : public Function< FunctionSpaceType, myRHSFunction<FunctionSpaceType> >
{
  // Implementation...
};
\end{lstlisting}

\begin{lstlisting}
template <class FunctionSpaceType>
class myExactSolution
    : public Function< FunctionSpaceType, myExactSolution<FunctionSpaceType> >
{
  // Implementation...
};
\end{lstlisting}
Note: These classes are derived from the \texttt{Function} class by using the Barton-Nackman-Trick in the second template parameter.

Now we put all problem data together in two classes:
\begin{lstlisting}
// right hand side and exact solution for poisson problem
template <class FunctionSpaceType>
struct PoissonProblem
{
    // type of right hand side
    typedef myRHSFunction<FunctionSpaceType> RHSFunctionType;

    // type of exact solution
    typedef myExactSolution<FunctionSpaceType> ExactSolutionType;
};

// for the L2-Projection use the exact solution also for the right hand side
template <class FunctionSpaceType>
struct L2ProjectionProblem
{
    // type of right hand side, use exact solution(!)
    typedef myExactSolution<FunctionSpaceType> RHSFunctionType;

    // type of exact solution
    typedef myExactSolution<FunctionSpaceType> ExactSolutionType;
};
\end{lstlisting}





\subsection{Right hand side assembler and boundary treatment in \texttt{boundary.hh} }
You can find more details about this topic in chapter 1 of this skript (in german): \url{http://www.mathematik.uni-freiburg.de/IAM/Teaching/ubungen/sci_com_SS06/skriptum.pdf}.

\subsubsection*{The function \texttt{assembleRHS()} }
This function assembles the right hand side for a given \texttt{function} and returns the result in \texttt{discreteFunction}.
The (given) right hand side function $f$ itself is implemented in \texttt{problem.hh}.
\begin{lstlisting}
template< int polOrd,
          class FunctionType,
          class DiscreteFunctionType>
static void assembleRHS( const FunctionType &function,
                         DiscreteFunctionType &discreteFunction )
{
// Implementation...
}
\end{lstlisting}

\subsubsection*{The function \texttt{boundaryTreatment()} } \label{boundaryTreatment}
Sets the dirichlet points to the corresponding values of \texttt{boundaryValue}. \texttt{boundaryValue} can be arbitrary, but  \texttt{langrange.cc} uses the (given) exact solution $u$ for setting the boundary values.
The exact solution $u$ itself is implemented in \texttt{problem.hh}.
\begin{lstlisting}
template< class EntityType,
          class FunctionType,
          class DiscreteFunctionType >
static void boundaryTreatment( const EntityType &entity,
                               const FunctionType& boundaryValue,
                               DiscreteFunctionType &rhs )
{
// Implementation...
}
\end{lstlisting}




\section{Implementing Laplace- and MassOperator}
Remember the following lines of the algorithm in \texttt{lagrange.cc}:
\begin{lstlisting}
// type of inverse operator
typedef OEMBICGSTABOp<DiscreteFunctionType,OperatorType> InverseOperatorType;

// assemble right hand side and so on ...
// ...

// create problem operator (Laplace or MassMatrix Operator)
OperatorType problemOperator( discreteFunctionSpace );

// apply inverse operator (solve the linear system)
InverseOperatorType cg( problemOperator, 1e-8, 1e-8, 20000, verbose );
cg( rhs, solution );
\end{lstlisting}
This means, for solving our problem with the inverse operator \texttt{cg}, we first have to provide a suitable problem operator.
How this can be done for the Poisson Problem is shown in the following subsections.


\subsection{An abstract matrix operator in \texttt{matrixoperator.hh} }
In this file the class \texttt{MatrixOperator} is implemented. This class is the base class for the problem operator classes \texttt{LaplaceOperator} and \texttt{MassOperator}, which are described below. The derived classes have to implement at least the (protected) method \texttt{assembleLocalMatrix()}.

\subsubsection{Private member variables}
\begin{itemize}
 \item \lstinline!const DiscreteFunctionSpaceType &discreteFunctionSpace_!: reference to the used discrete function space.

 \item \lstinline!mutable MatrixObjectImp matrixObj_!: pointer to the system matrix.

 \item \lstinline!mutable bool matrix_assembled_!: flag indicating whether the system matrix has already been assembled.

 \item \lstinline!const bool boundaryCorrect_!: true if boundary correction should be done.
\end{itemize}


\subsubsection{The constructor \texttt{MatrixOperator::MatrixOperator()} }
The constructor of this class is protected, so you cannot create any instances of this class.
\begin{lstlisting}
MatrixOperator( const DiscreteFunctionSpaceType &discreteFunctionSpace, bool bnd )
    : discreteFunctionSpace_( discreteFunctionSpace ),
      matrixObj_( discreteFunctionSpace, discreteFunctionSpace , "" ),
      matrix_assembled_( false ),
      boundaryCorrect_(bnd)
{
}
\end{lstlisting}


\subsubsection{Public methods}
The public methods are all very easy.
\begin{itemize}
 \item \texttt{MatrixOperator::operator() ()}\\
       \texttt{MatrixOperator::multOEM()}\\
    Apply the operator: Multiplies a vector with the system matrix and returns the result.

 \item \texttt{MatrixOperator::systemMatrix()}\\
    Returns the system matrix \texttt{matrixObj\_}. Calls \texttt{MatrixOperator::assemble()} before, if needed.

 \item \texttt{MatrixOperator::preconditionMatrix()}\\
    Returns a reference to the preconditioning matrix.

 \item \texttt{MatrixOperator::print()}\\
    Prints the system matrix into a stream.

 \item \texttt{MatrixOperator::discreteFunctionSpace()}\\
    Returns the used discreteFunctionSpace.

 \item \texttt{MatrixOperator::assemble()}\\
    Allocates memory for the system matrix \texttt{matrixObj\_} and then
    calls the method \texttt{MatrixOperator::assembleOnGrid()}.
\end{itemize}


\subsubsection{Protected methods}
\begin{itemize}
 \item \texttt{MatrixOperator::assembleOnGrid()}\\
    Runs over all entities, calls on each entity \texttt{MatrixOperator::assembleOnEntity()} and calls,
    if needed, \texttt{MatrixOperator::boundaryCorrectOnEntity()}.
    Furthermore, prints out the time needed for the hole matrix to assemble.

 \item \texttt{MatrixOperator::assembleOnEntity()}\\
    Performs matrix assemble for one entity. Creates for that entity a local matrix and calls
    \texttt{MatrixOperator::assembleLocalMatrix()} on the entity.

 \item \texttt{MatrixOperator::assembleLocalMatrix()}\\
    Uses the Barton-Nackman Trick to avoid using virtual functions: Just calls the corresponding method
    \texttt{assembleLocalMatrix()} in the derived class. The derived class has to implement this method!

 \item \texttt{MatrixOperator::boundaryCorrectOnEntity()}\\
    Modifies the system matrix in order to implement the dirichlet boundary condition. Compare with the function
    \texttt{boundaryTreatment()} on page \pageref{boundaryTreatment}, where the corresponding thing was done
    for the right hand side.
\end{itemize}


\subsubsection{Private methods}
\begin{itemize}
 \item \texttt{MatrixOperatorType\& asImp()}\\
       \texttt{const MatrixOperatorType\& asImp() const}\\
    Performs the Barton-Nackman Trick to avoid virtual functions. Looks strange, but works, so don't worry about!
\end{itemize}




\subsection{The Laplace operator in \texttt{laplace.hh} }
The struct \texttt{LaplaceOperatorTraits} at the beginning of the file contains only some typedefs,
so let's start directly with the  class \texttt{LaplaceOperator}.
\begin{lstlisting}
template< class DiscreteFunctionImp >
class LaplaceOperator : public MatrixOperator< LaplaceOperatorTraits<DiscreteFunctionImp> >
\end{lstlisting}
This class is derived from the class \texttt{MatrixOperator} (see \texttt{matrixoperator.hh}) to implement the Laplace operator.
The typedefs at the beginning of the class are not so important.

Most of the work is already be done in the base class, here we have only to implement the method \texttt{LaplaceOperator::assembleLocalMatrix()}. For details about that method we refer again to
this skript: \url{http://www.mathematik.uni-freiburg.de/IAM/Teaching/ubungen/sci_com_SS06/skriptum.pdf}, chapter 1.3.2.



\subsection{The mass matrix operator in \texttt{massmatrix.hh} }

This file holds another class derived from \texttt{MatrixOperator}, but this time to implement the mass matrix operator.
The hole file is very similar with the file \texttt{laplace.hh}, only the method \texttt{MassOperator::assembleLocalMatrix()}
is (of course) different.

















  %Chapter 4
  \chapter{An LDG solver for Advection-Diffusion Equations}
  \newcommand{\velo}{\mathbf{a}}
\newcommand{\timedomain}{\mathbf{T}}
This chapter is to introduce the \texttt{Pass} concept in \Fem\ and to explain
the steps that are necessary in order to implement a Local Discontinuous
Galerkin solver using the \texttt{LocalDGPass} class. In the end it is
shown how the solution can be analysed visually with help of the
\texttt{DataWriter} class.

\section{Advection-Diffusion Equation}

The example problem is a scalar advection-diffusion equation on $\Omega =
[0,1]^3$
\begin{equation}
  \label{eq:advectdiff}
  \begin{aligned}
    \partial_{t} u + \nabla \cdot (\velo u) - \varepsilon \Delta u 
    {}          &=  0            & \mbox{in}&\ \Omega \times \timedomain \\
    u           &=  g_{D}        & \mbox{on}&\ \partial\Omega \times \timedomain \\
    u(0,\cdot)  &=  u_{0}        & \mbox{in}&\ \Omega \times \{0\},
  \end{aligned}
\end{equation}
with $u$ living in a function space $V$ and $\velo:=(0.8,0.8,0)^t$. An exact
solution is specified by
% definitions for equation
\newcommand{\ccos}{\tilde{c}}
\newcommand{\csin}{\hat{c}}
\begin{equation}
  u(x,t) = \sum_{i=1}^{2} \exp(-\varepsilon t \pi^2 |{c^{i}|})
  \left(\prod_{j=1}^{3} \ccos_{j}^{i}\cos(c_{j}^{i}\pi(x_{j}-\velo_{j} t)) + 
  {}              \csin_{j}^{i}\sin(c_{j}^{i}\pi(x_{j}-\velo_{j} t)) \right),
  \label{eq:exactsolution}
\end{equation}
where
\begin{align}
  c^{1}    &:=(2,1,1.3)^{t}      & c^{2}&:=(0.7,0.5,0.1) \\
  \csin^{1}&:=(0.8,0.4,-0.4)^{t} & \csin^{2} &:= (0.2,0.1,0.2) \\
  \ccos^{1}&:=(0.6,1.2,0.1)^{t}  & \ccos^{2} &:= (0.9,0.3,-0.3).
  \label{eq:constants}
\end{align}
Therefore the initial and boundary functions are defined by
\begin{align}
  u_{0}(x) &= u(x,0) & g_{D}(x,t) = u|_{\partial\Omega}(x,t).
  \label{eq:initialboundary}
\end{align}
Note that the problem is also implemented for $2$ dimension by projecting
everything onto the first two coordinates. 
The discretization of the problem is done as described in
\cite[ch.4]{wissrechnen}.

The LDG Ansatz uses auxiliary functions $u_{0}, u_{2} \in V$ and $u_{1} \in V^{3}$
%
\begin{equation}
  \label{eq:ldgansatz}
  \begin{aligned}
    u_{0} &= u \\
    u_{1} &= - \sqrt{\varepsilon} \nabla u_{0}  \\
    u_{2} &= - \nabla \cdot (\velo u_{0} + \sqrt{\varepsilon} u_{1})  \\
    \partial_{t} u &= u_{2}.
  \end{aligned}
\end{equation}
%
Now $u_{0}$ and $u_{1}$ can be projected onto the Discontinous Galerkin
function spaces $V_{h} := \{\varphi_{1},\dots \varphi_{n}\}$ resp.
$[V_{h}]^{3}=\{\psi_{1},\dots\ \psi_{n^{3}}\}$ with discrete domain in space. The equations
in (\ref{eq:ldgansatz}) can then be rewritten as a variational formulation 
%
\newcommand{\ftilde}{\tilde{f}}
\begin{align}
  \label{eq:variational1}
  \int_{T}^{} u_{1} \psi &= \int_{T}^{}\underbrace{\sqrt{\varepsilon}u_{0}}_{f_{1}(u_{0})} \nabla \cdot \psi
  {}                       - \int_{\partial T}^{}\underbrace{\sqrt{\varepsilon}u_{0}}_{\ftilde_{1}(u_{0})} \psi \cdot n 
  {}                                   && \forall{\psi}\in [V_{h}]^{3}\\
  \label{eq:variational2}
  \int_{T}^{} u_{2} \varphi &= \int_{T}^{} \underbrace{\left(\velo u_{0} + \sqrt{\varepsilon} u_{1}\right)}_{
  {}                                                      f_{2}(u_{0},u_{1})} \cdot \nabla \varphi
  {}                         - \int_{\partial T}^{} \underbrace{\left(\velo u_{0} + \sqrt{\varepsilon} u_{1}\right)}_{
  {}                                                      \ftilde_{2}(u_{0},u_{1})} \varphi \cdot n
  {}                                         && \forall{\varphi} \in V_{h}
\end{align}
for every grid cell $T$. 
%
To complete the discretization in the space domain, the numerical fluxes
between the cell interfaces must be defined to approximate $\ftilde_{1}$ and
$\ftilde_{2}$ defined in (\ref{eq:variational1}) and (\ref{eq:variational2}).
In this example we choose
%
\begin{align}
  \label{eq:flux1}
  \ftilde_{1,h}(u_{0}) &= 
  \left\{
  \begin{array}{cl}
    \sqrt{\varepsilon} g_{D}(x)  & \mbox{on}\ \partial\Omega \\
    \sqrt{\varepsilon} \{u_{0}\} & \mbox{else}
  \end{array}
  \right.
\\
  \label{eq:flux2}
  \ftilde_{2,h}(u_{0},u_{1}) &=
  \left\{
  \begin{array}[]{cl}
    \velo\, g_{D}(x) + \sqrt{\varepsilon}u_{1}(x^{-}) & \mbox{on}\ \partial\Omega \\
    w(\velo, u_{0})  + \sqrt{\varepsilon}\{u_{1}\} & \mbox{else}
  \end{array}
  \right.
\end{align}
%
where
%
\begin{equation}
  \{u\} := \frac{1}{2}(u(x^{+}) + u(x^{-})) \quad\mbox{and}\quad
  w(\velo, u_{0}) := \left\{
  \begin{array}[]{cl}
    \velo u_{0}(x^{+}) & \mbox{if}\ \velo \cdot n \leq 0 \\
    \velo u_{0}(x^{-}) & \mbox{if}\ \velo \cdot n > 0
  \end{array}
  \right.
  \label{eq:meanupwind}
\end{equation}
%
define a mean respectively an upwind function over the cell interfaces.
%
Now an operator $L : V_{h} \to V_{h}$ can be defined as a concatenation
$L:=\Pi \circ L_{2}\circ L_{1}$ of the two operators 
%
\begin{equation}
  \label{eq:passes}
  \begin{aligned}
    L_{1}:& V_{h} \to V_{h} \times V_{h}^{3} && (u_{0}) \mapsto (u_{0}, u_{1}) \\
    L_{2}:& V_{h} \times V_{h}^{3} \to V_{h} \times V_{h}^{3} \times V_{h} && (u_{0}, u_{1}) \mapsto (u_{0}, u_{1}, u_{2})
  \end{aligned}
\end{equation}
%
and the projection operator $\Pi$ that projects to the last component. This allows to
shorten the problem formulation to
%
\begin{equation}
  \partial_{t} u + L[u] = 0.
  \label{eq:advectdiffshort}
\end{equation}
%
This is an ODE in the time domain and can be solved with standard ODE solvers
requiring evaluations of the discrete operator $L$ in each time step. The
so-called ``methods of lines'' is further examined in
\cite[ch.4.3]{wissrechnen}.  
%
\section{Implementation overview}\label{se:overview}

The source code of the implmentation described in this section is located in
the directory \lstinline!dune-femhowto/src\_localdg!. It consists of $7$ source
files:
%
\begin{itemize}
  \item \lstinline!problem.hh! contains the class \lstinline!U0! which defines
    the exact solution $u$ as defined in (\ref{eq:exactsolution}) and $g_D, u_0$
    as defined in (\ref{eq:initialboundary}). Those can be called through the
    following methods:
      \lstinputlisting[numbers=left,numberstyle=\tiny,numbersep=5pt,linerange={94-94,111-111}]{../src_localdg/problem.hh}
    
  \item \lstinline!models.hh! contains the class
    \lstinline!AdvectionDiffusionModel! that describes the model data given in
    equation (\ref{eq:ldgansatz}) and the class \lstinline!UpwindFlux! which is
    just a helper class for the upwind flux given in (\ref{eq:meanupwind}).

  \item \lstinline!discretemodels.hh! contains two classes
    \lstinline!AdvDiffDModel1! and \lstinline!AdvDiffDModel2! that describe the
    two passes of the LDG implementation, i.e. the terms and fluxes given by
    (\ref{eq:variational1}--\ref{eq:flux2}). For further details to these classes
    see section \ref{se:localdg}.

  \item \lstinline!advectdiff.hh! provides the LDG Operator
    \lstinline!DGAdvectionDiffusionOperator! that is constructed out of the
    models given in \lstinline!discretemodels.hh!. For further details, refer
    to section \ref{se:localdg}.

  \item \lstinline!dgtest.cc! is the source file for the main program and is
    explained in the next section.

  \item \lstinline!datadisp.cc! and \lstinline!stuff.cc! are functions with helper class
    for pretty printing of EOC Data respectively data visualisation. Further
    details on this topic can be found in section \ref{se:outputvisualisation}.
%
\end{itemize}

\section{Main Loop} \label{se:mainloop}
%
The major typedefs for problem and model definitions are included via the
\lstinline!models.hh! header file. If the program shall run with a different
ODE solver e.g., the definitions that can be seen in the header files excerpt
in listing \ref{lst:modelshh} need to be adapted.
%
\begin{lst}[End of ../src\_localdg/models.hh] \label{lst:modelshh}\mbox{}
  \lstinputlisting[numbers=left,numberstyle=\tiny,numbersep=5pt,linerange={325-348}]{../src_localdg/models.hh}
\end{lst}
%
The main loop can be found in the file \lstinline!dgtest.cc!. An explanation
of the most important lines in this file, follows after listing
\ref{lst:dgtestcc}.
%
\begin{lst}[../src\_localdg/dgtest.cc] \label{lst:dgtestcc}\mbox{}
  \lstinputlisting[numbers=left,numberstyle=\tiny,numbersep=5pt]{../src_localdg/dgtest.cc}
\end{lst}
%
In lines \ref{dg:param0}-\ref{dg:param1} runtime options specified in a
parameter file are read in with help of the \lstinline!Dune::Parameter! class.
The default name for the parameter file is ``parameter'' (as defined in line
\ref{dg:paramfile}) and can be changed through the command line.

Lines \ref{dg:cfl0}-\ref{dg:cfl1} set the CFL constant that is needed to
preserve strong stability of the ODE solver. This constant can also be
specified explicitly in the parameterfile. It is known by the TimeProvider
\lstinline!tp! that in turn is given together with the discrete operator for
the space domain \lstinline!dg!  to the ODE solver in lines \ref{dg:dgop} and
\ref{dg:timeprovider}. Each time the ODE solver calls one of its methods
\lstinline!initialize! or \lstinline!solve! (see lines \ref{dg:odeinit},
\ref{dg:odesolve}), it evaluates the operator \lstinline!dg! and afterwards
asks the operator for a new time step estimate by calling its method
\lstinline!provideTimeStepEstimate!. This estimate multiplied by the CFL
constant then turns out to be the next timestep length. Note that in lines
\ref{dg:maxtimestep1} and \ref{dg:maxtimestep2} the maximum for this timestep
length is bounded by the value in \lstinline!maxTimeStep!, which is controlled
via the parameterfile.

For visualisation of the found solution the program writes the time snapshots
to the harddisk. This snapshot data can then be postprocessed by programs like
GRAPE or Paraview.  The user hast the ability to change the output directory
for the snapshots, the interval at which the snapshots should be generated and
the output format in the parameterfile. The options are all handled by an
instance of the \Fem\ class \lstinline!DataWriter! that is created in lines
\ref{dg:datawriterdef0}-\ref{dg:datawriterdef1} and requested to write
snapshots on the harddisk in lines \ref{dg:datawrite1} and \ref{dg:datawrite2}. 

\section{Setting up an LDGPass} \label{se:localdg}

This section is to shed light on the construction of the operator
\lstinline!dg! that obviously is the heart of this LDG solver.

In order to set up the operator $L$ as it is defined in (\ref{eq:passes})
the \Fem\ Pass concept is used. 

The passes $L_{1}$ and $L_{2}$ are implemented as \texttt{LocalDGPass}
instances. A \texttt{LocalDGPass} solves an equation of the form
%
\begin{align} v + \text{div}(f(x,u)) + A(x,u)\nabla u &= S(x,u) &&\mbox{in}\
  \Omega.  \label{eq:LocalDGPass} \end{align}
%
with the argument $u$ and the computed solution $v$. Both required passes (see
(\ref{eq:ldgansatz})) are in this form. A look at the weak formulation of
equation (\ref{eq:LocalDGPass}) reveals what methods a model which is passed to
the constructor of a \texttt{LocalDGPass} instance needs to specify:
%
\begin{equation} \label{eq:fourmethods} \int_{T}^{} v \varphi = -
  \int_{\partial T}^{} \underbrace{\left( f(x,u) + A(x,u) \cdot n
  \right)}_{\substack{\mathtt{numericalFlux}\\\mathtt{boundaryFlux}}} \varphi +
  \int_{T}^{} \overbrace{f(x,u)}^{\mathclap{\mathtt{analyticalFlux}}} \cdot
  \nabla \varphi + \int_{T}^{} \underbrace{\left(S -
  A(x,u)\right)}_{\mathtt{source}} \nabla u \varphi.  \end{equation}
%
The header file \lstinline!advectdiff.hh! defines the two models that implement
these four methods \lstinline!analyticalFlux!, \lstinline!source!,
\lstinline!numericalFlux! and \lstinline!boundaryFlux!. Each of these methods
describes the discrete version of its corresponding term in the variational
formulation (\ref{eq:fourmethods}). The second pass of this implementation e.g.
has no \lstinline!source! term, the \lstinline!analyticalFlux! represents
$f_{2}(u_{0}, u_{1})$ (see (\ref{eq:variational2})) and the
\lstinline!numericalFlux! resp. the \lstinline!boundaryFlux! together represent
the flux function $f_{2,h}(u_{0},u_{1})$ (see (\ref{eq:flux2})), where the
first one implements the flux on interfaces between inner cells and the second
one on the boundary domain. Listing \ref{lst:advdiffd2} shows the actual
implementation of this model.
%
\begin{lst}[AdvDiffDModel2 in ../src\_localdg/discretemodels.hh]
  \label{lst:advdiffd2}\mbox{}
  \lstinputlisting[numbers=left,numberstyle=\tiny,numbersep=5pt,linerange={254-403}]{../src_localdg/discretemodels.hh}
\end{lst}
%
Note that in line \ref{dm:hasSource} the model notifies the
\lstinline!LocalDGPass! that there is no source term, and therefore the
\lstinline!source! method is never called and does not need to be implemented.
All the implemented fluxes have a parameter of type \lstinline!ArgumentTuple!.
This is the argument to the flux function which for the second pass lives in
the product of two function vector spaces. In order to access the correct
component of the argument, there is an elegant solution using the
\lstinline!Int2Type! helper class: Every Pass gets a unique id when it is
defined (see line \ref{ad:passids}) and the \lstinline!Int2Type! class can
convert this unique id to a unique type as done in lines \ref{dm:int2type0} and
\ref{dm:int2type1}. These types then allow \lstinline!ArgumentTuple!s to be
used like arrays, as in line \ref{dm:int2typeusage} e.g., where
\lstinline!uLeft[u0Var]! returns $u_{0}(x^{-})$ and \lstinline!uLeft[u1Var]!
returns $u_{1}(x^{-})$. Alternatively there is a \lstinline!Selector! class in
\Fem\ which also provides accessors to any component of a tuple but which is
more cumbersome.

With models defined for both passes the implementation of the LDG operator
becomes quite simple.
%
\begin{lst}[../src\_localdg/advectdiff.hh] \label{lst:advectdiffhh}\mbox{}
  \lstinputlisting[numbers=left,numberstyle=\tiny,numbersep=5pt]{../src_localdg/advectdiff.hh}
\end{lst}
%
Lines \ref{ad:typedefpass1} and \ref{ad:typedefpass2} contain the typedefs for
both passes specifying the underlying model and the pass id of the
\lstinline!LocalDG! pass. The initialisation of the passes in lines
\ref{ad:initialisepass1} and \ref{ad:initialisepass2} is equally
straightforward.  Line \ref{ad:callpass2} shows how a pass can be evaluated
like an \lstinline!Operator! calling all previous passes.
%
%If you are interested in writing your own Dune Pass you should refer to
%section \ref{se:pass_details} where the implementation of the \lstinline!Pass!
%and \lstinline!LocalDGPass! classes is explained in more detail.
%
%\section{Pass details} \label{se:pass_details}

\section{Visualisation and EOC Output} \label{se:outputvisualisation}

The source code shipped with this document contains a file called
\lstinline!datadisp.cc! which implements a GRAPE interface for the data
produced by the \lstinline!dgtest! program. It also provides a file
\lstinline!manager.replay! which preconfigures GRAPE for watching the snapshots
as an animation. Just click on the play button.
%
\begin{lstlisting} 
  ./datadisp paramfile:parameter 0 25 
\end{lstlisting}
%
starts the visualisation, where the last number must be less than or equal to
the number of snapshots produced by \lstinline!dgtest!.

Furthermore, \lstinline!dgtest! produces on every run, a \TeX file called
\lstinline!eoc.tex!, which comprises statistics about the numerical convergence
of the implemented LDG scheme.

% vim:set ft=tex:


  %Chapter 5
  \chapter{Documentation and reference guide for \Fem}
  The complete documentation of \Fem\ is done using the tool \texttt{doxygen}.
You can find this documentation online under \url{http://www.mathematik.uni-freiburg.de/IAM/Research/projectskr/dune/doc/html/index.html}.
You can also build your own local \texttt{doxygen} documentation, by removing the \texttt{--disable-documentation} command in your \texttt{config.opts} file (see listing on page \pageref{config.opts})

Some hints about using the \texttt{doxygen} documentation:
\begin{itemize}
\item The best way to start is from the page Modules (\url{http://www.mathematik.uni-freiburg.de/IAM/Research/projectskr/dune/doc/html/modules.html}) which gives you access to the documentation by category.

\item A list of the central interface classes can be found here: \url{http://www.mathematik.uni-freiburg.de/IAM/Research/projectskr/dune/doc/html/interfaceclass.html}.

\item A summary of the main features and concepts can be found here: \url{http://www.mathematik.uni-freiburg.de/IAM/Research/projectskr/dune/doc/html/group__FEM.html}.

\item Newly added implementations are linked on this page: \url{http://www.mathematik.uni-freiburg.de/IAM/Research/projectskr/dune/doc/html/newimplementation.html}.

\item Some remarks about writting dune-fem documentation is found here: \url{http://www.mathematik.uni-freiburg.de/IAM/Research/projectskr/dune/doc/html/DocRules.html}.

\item And finally some notes on using subversion (svn):
\url{http://www.mathematik.uni-freiburg.de/IAM/Research/projectskr/dune/doc/html/svnhelp.html}
\end{itemize}


  %Still very incomplete, use doxygen documentation instead!
  %\chapter{Introduction}

  The implementation of a discrete model of a partial differential equation in \Fem\ is based on the
  general concept of
    \begin{itemize}
    \item function spaces, 
    \item functions and
    \item operators, that act on functions.
  \end{itemize}

%%%%%%%%%%%%%%%%%%%%%%%%%%%%%%%%%%%%%%%%%%%%%%%%%%%%%%%%%%%%%%%%

\chapter{Functions and FunctionSpaces}

%%%%%%%%%%%%%%%%%%%%%%%%%%%%%%%%%%%%%%%%%%%%%%%%%%%%%%%%%%%%%%%%

\section{Abstract definition of function spaces and functions}

  A function space $V$ in our concept is a set of mappings from the domain $D:=\mathbb{K}_D^d$ to the range $R:=\mathbb{K}_R^n$, e.g.
  \begin{eqnarray}
    V := \{ u : \mathbb{K}_D^d \to \mathbb{K}_R^n \} \nonumber 
  \end{eqnarray}
  Here, $\mathbb{K}_D$ denotes the domain field, $\mathbb{K}_R$ the range field and  $d,n$ the dimensions of the domain and range, respectively. To further specify the function space, additional properties can be added, e.g. the function are in $C^m$ or do belong to the Sobolev Space $H^m$. \index{function space} \index{domain} \index{range} \index{domain field} \index{range field}

  A discrete function space \index{discrete function space} $V_h$ with finite dimension $m$ is a subset of a function space with the property that the functions are defined locally on the elements $e$ of the underlying computational grid $\cal T$. If $\hat e $ denotes the reference element of $e$ and $F_e$ the mapping $F_e: \hat e \to e$ (for further details refer to \citep{GitterPapier:06}), we define the local base function set $V_{\hat e}$ on the reference element $\hat e$ through
  \begin{eqnarray}
    V_{\hat e} := \{\varphi_1, \ldots, \varphi_{{\rm dim}(V_{\hat e})}\}. \nonumber 
  \end{eqnarray}
  The discrete function space $V_h$ is then given as 
\begin{eqnarray}
    V_h := \Big\{ u_h \in V: u_h|_e := u_e := \sum_{\varphi \in V_{\hat e}} g(u_{e,\varphi}) \ \varphi \circ F_e^{-1}, \ {\rm for \ all } \ e \in {\cal T}\big\}. \nonumber 
\end{eqnarray}
  We call $V_e := {\rm span}\{\varphi \circ F_e^{-1}: \varphi \in V_{\hat e}\}$ a local function space, $u_e \in V_e$ a local function, and ${\rm DOF}_e := \{u_{e,\varphi}, \varphi \in V_{\hat e}\}$ the set of local degrees of freedom. In oder to incorporate global properties of the discrete function space, the function space has to provide a mapping $g$ between the local degrees of freedom (${\rm DOF}_e$) and the global degrees of freedom ${\rm DOF} := \{ u_i: i = 0, \cdots, m\}$. \index{function space!local} \index{DOF, degree of freedom}

  We summarize, that a discrete function space $V_h$ is determined by a function space $V$, a grid $\cal T$, the base function sets $V_{\hat e}$ for all reference elements $\hat e$ and the mapping $g$ from local to global degrees of freedom. A discrete function $u_h \in V_h$ is accordingly defined as a set of local functions $u_e$ where a local function provides access to the local degrees of freedom (${\rm DOF}_e$).


  \section{FunctionSpaces and Mappings}

   A function space $V$ in our concept is a set of mappings from the domain $D:=\mathbb{K}_D^d$ to the range $R:=\mathbb{K}_R^n$, e.g.

  \begin{eqnarray}
    V := \{ u : \mathbb{K}_D^d \to \mathbb{K}_R^n \} \nonumber 
  \end{eqnarray}

  Here, $\mathbb{K}_D$ denotes the domain field, $\mathbb{K}_R$ the range field and  $d,n$ the dimensions of the domain and range, respectively. 


  \section{Grids and GridParts}

  In this section we will remember how to create grids and GridParts. The abstract definition of \Dune\ grids can be found  in \citep{GitterPapier:06}. 

  Remember, a hierarchic grid $\Th$ with $L+1$ levels consists of $L+1$ grids 
  $$\Th := \left\{\mathcal{T}_0, \mathcal{T}_1, \ldots \mathcal{T}_L\right\},$$
  and $\mathcal{T}_i$ is formed up by the codim 0 entities of level $i$, $i=0,\ldots,L$. 

  We call $\mathcal{T}_0$ the macro grid \index{grid!macro}. Another important grid consisting of codim 0 entities of a \Dune\ grid is the leaf grid \index{grid!leaf} $\mathcal{T}_{leaf}$. It is made up by the codim 0 entities that have no children.  

  Each of these grids, $\mathcal{T}_0, \ldots, \mathcal{T}_L, \mathcal{T}_{leaf}$, is suitable for computations within the DUNE FEM context! With the GridPart class you can choose one these grids. This is necessary for the definition of a function space. Imagine a GridPart as a view on a part of the underlying grid. Most of the methods of the GridPart class then replace the methods offered by the grid e.g. when working with the discrete functions. The interface methods are
 

\section{FunctionSpaces}

  In this section we will show all the defintions needed to define a function space and a function and how to evaluate function.

%%%%%%%%%%%%%%%%%%%%%%%%%%%%%%%%%%%%%%%%%%%%%%%%%%%%%%%%%%%%%%%%

\chapter{DiscreteFunctions and Spaces}

%%%%%%%%%%%%%%%%%%%%%%%%%%%%%%%%%%%%%%%%%%%%%%%%%%%%%%%%%%%%%%%%


%%%%%%%%%%%%%%%%%%%%%%%%%%%%%%%%%%%%%%%%%%%%%%%%%%%%%%%%%%%%%%%%

\chapter{Operators}

%%%%%%%%%%%%%%%%%%%%%%%%%%%%%%%%%%%%%%%%%%%%%%%%%%%%%%%%%%%%%%%%


%%%%%%%%%%%%%%%%%%%%%%%%%%%%%%%%%%%%%%%%%%%%%%%%%%%%%%%%%%%%%%%%

\chapter{Passes}

%%%%%%%%%%%%%%%%%%%%%%%%%%%%%%%%%%%%%%%%%%%%%%%%%%%%%%%%%%%%%%%%

\section{General concept}

\section{An example}


%%%%%%%%%%%%%%%%%%%%%%%%%%%%%%%%%%%%%%%%%%%%%%%%%%%%%%%%%%%%%%%%

\chapter{Numerical methods}

%%%%%%%%%%%%%%%%%%%%%%%%%%%%%%%%%%%%%%%%%%%%%%%%%%%%%%%%%%%%%%%%

\section{Laplace equation, implicit computation}

\section{The finite volume method}


\section{The LDG method}


%%%%%%%%%%%%%%%%%%%%%%%%%%%%%%%%%%%%%%%%%%%%%%%%%%%%%%%%%%%%%%%%

\chapter{Solvers}

%%%%%%%%%%%%%%%%%%%%%%%%%%%%%%%%%%%%%%%%%%%%%%%%%%%%%%%%%%%%%%%%

\section{Forward Euler}


\section{CD method}


\section{IMEX method}





  %I think there is no need anymore for this chapter. In my opinion, the topic is already done in the chapters 2 and 3.
  %%%%%%%%%%%%%%%%%%%%%%%%%%%%%%%%%%%%%%%%%%%%%%%%%%%%%%%%%%%%%%%%%

\chapter{Problem implementation}

%%%%%%%%%%%%%%%%%%%%%%%%%%%%%%%%%%%%%%%%%%%%%%%%%%%%%%%%%%%%%%%%

In this part we want to show how to work with \Dune\ \Fem. Use it as a introduction and step by step tutorial to \Fem.

The \Fem\ interface was designed to translate mathematical formulations easily to computable algorithms. We will discuss several examples of classical mathematical problems like the poisson equation and how to implement different numerical schemes. 

As mentioned abouve we will discuss the poisson equation in the first chapter of this quick guide. This equation can be treated with both finite element and finite volume methods. Another possibility is to use the Local Discontinuous Galerkin method.

\Fem\ includes solvers for linear problems and for the solution of ODEs. They are discussed in chapter 7.

At the end of this part we will look at the Input/Output methods of \Fem\ and how to visualize your data with Grape\index{Grape}.


\section{Getting started}
  \begin{lst}[File ../sources/gettingstarted.cc] \mbox{}
    \lstinputlisting[basicstyle=\ttfamily\scriptsize,numbers=left, numberstyle=\tiny, numbersep=5pt]{../sources/gettingstarted.cc}
  \end{lst}
  In our first example, we will create a grid and a discrete function space that is defined on the grid. This program is very simple - well, it doesn't do anything. 

  So, why all the typedefs? Remember, \Dune\ was designed to provide efficient implementations using generic programming techniques. Therefore, FunctionSpaces, Functions, Operators, etc. are represented by classes with several template arguments. In order to construct one these objects, you have to declare all the types needed in the beginning. 

  \Fem\ helps you managing types. You can acces the most important types everywhere as they are public on every class. Use the types offered there and you will avoid mistakes and have less problems.


\section{An application}
In this section we present a simple application of some of the dune-fem basic concepts.We want to calcuate:
\begin{equation}
\int_\Omega \nabla\cdot F dx 
\end{equation}\label{div}
for a given discrete Function $F:\Omega\rightarrow\mathbb{R}^n$.
This will be done in two ways:\\
First we calculate the Volume Interal in an element by element manner :\\
\begin{equation}
\int_\Omega \nabla\cdot F dx =\sum_{e\in\mathcal{E}^0_{leaf}} \int_e \nabla\cdot F dx
\end{equation}\\
Then we use the divergence theorem:$\int_\Omega \nabla\cdot F dx= \int_{\partial\Omega}F\cdot \nu ds$ to write \ref{div}as a boundary integral: $\int_\Omega \nabla\cdot F dx =\sum_{e\in\mathcal{E}^1_{bnd}} \int_e  F\cdot\nu ds$
  
  
  
  
  
  







  %%%%%%%%%%%%%%%%%%%%%%%%%%%%%%%%%%%%%%%%%%%%%%%%%%%%%%%%%%%%%%%%

  % Appendix, bibliography, index, etc.

  %%%%%%%%%%%%%%%%%%%%%%%%%%%%%%%%%%%%%%%%%%%%%%%%%%%%%%%%%%%%%%%%

  \bibliographystyle{plain}
  \bibliography{bibliography}

  \bigskip
  For a more complete list of publications see also: \url{http://www.dune-project.org/publications.html}

%  \printindex

\end{document}
